%% =============================================================================
%% manuscript.tex — Expositiones Mathematicae Short Research Note
%% Title:    The continuous-coefficient Jensen equation
%% Subtitle: A note on three vestigial regularity hypotheses
%% Author:   [AUTHOR]
%% Source:   satellites/o3-expositiones/01-manuscript.md
%% Date:     2026-06-07
%%
%% Compilation:
%%   pdflatex manuscript
%%   bibtex   manuscript
%%   pdflatex manuscript
%%   pdflatex manuscript
%%
%% Class:    Elsevier's elsarticle class. Two modes:
%%   * `review,12pt` for the SUBMISSION pdf (double-spaced, line-numbered, single-column)
%%   * `preprint,12pt` for an author-preferred reading pdf (single-column, no line numbers)
%%   * `1p,12pt`     for the typeset look (publisher-style single column)
%%   * `5p,times`    for the typeset look (publisher-style 2-column)
%% The current setting is `preprint` which produces a clean reading copy.
%% For EES dispatch, switch to `review` (uncomment the corresponding line).
%% =============================================================================

% --- Submission mode (uncomment for the SUBMISSION pdf) ---
\documentclass[review,12pt,a4paper]{elsarticle}

% --- Reading mode (active by default) ---
% \documentclass[preprint,12pt,a4paper]{elsarticle}

% --- Final-typeset mode (uncomment for what the published article looks like) ---
% \documentclass[1p,12pt,a4paper,times]{elsarticle}

% ============= Packages =============
\usepackage[utf8]{inputenc}
\usepackage[T1]{fontenc}
\usepackage{lmodern}   % scalable Latin Modern fonts (required for microtype expansion)
\usepackage{amssymb,amsmath,amsthm}
\usepackage{microtype}
\usepackage{booktabs}
\usepackage{array}
\usepackage{tabularx}
\usepackage[hidelinks]{hyperref}

% ============= Theorem environments =============
\newtheorem{theorem}{Theorem}[section]
\newtheorem{corollary}[theorem]{Corollary}
\newtheorem{proposition}[theorem]{Proposition}
\newtheorem{lemma}[theorem]{Lemma}
\theoremstyle{remark}
\newtheorem{remark}[theorem]{Remark}

% ============= Math shortcuts =============
\newcommand{\R}{\mathbb R}
\newcommand{\Q}{\mathbb Q}
\newcommand{\N}{\mathbb N}
\newcommand{\E}{\mathbb E}

% ============= Front matter =============
\journal{Expositiones Mathematicae}

\begin{document}

\begin{frontmatter}

\title{The continuous-coefficient Jensen equation:\\
       A note on three vestigial regularity hypotheses}

\author[inst1]{[AUTHOR]\corref{cor1}}
\ead{[author@institution.country]}
\cortext[cor1]{Corresponding author.}

\affiliation[inst1]{organization={[INSTITUTION]},
                    addressline={[ADDRESS LINE 1]},
                    city={[CITY]},
                    postcode={[POSTAL CODE]},
                    country={[COUNTRY]}}

\begin{abstract}
The continuous-coefficient form of Jensen's functional equation on a real
interval --- the equation $p\,G(u_1)+(1-p)\,G(u_2)=G(p\,u_1+(1-p)\,u_2)$
imposed for every $u_1,u_2$ and every $p\in[0,1]$ --- has a one-line proof
of affineness via a single substitution that pins $G$ along the chord
through the endpoints, with no regularity hypothesis on $G$. This
contrasts sharply with the classical theory of Cauchy's additive
equation, where without a regularity hypothesis the Hamel-basis pathology
supplies non-affine solutions: the Polish school of the 1920s
(Sierpi\'nski, Steinhaus, Ostrowski) classifies the regularity hypotheses
(continuity, measurability, boundedness on a positive-measure set) that
recover affineness. We give the substitution explicitly, articulate
the dictionary of three classically required regularity hypotheses that
become vestigial under the continuous-coefficient form, and identify the
structural mechanism: the Hamel pathology lives at irrational coefficients
$p$, precisely where the continuous-coefficient equation forecloses on it.
We close by isolating the recurrence pattern of the trap --- Jensen's
inequality pushed to saturation --- and surveying three example application
areas (expected-utility theory, Shannon entropy's axiomatic
characterization, surrogate calibration) where the substitution is
the natural closure of the equation.
\end{abstract}

\begin{highlights}
\item The continuous Jensen equation $(\star)$ has a one-line proof: an endpoint substitution.
\item Three regularity hypotheses on $G$ become vestigial under continuous coefficients.
\item The Hamel-basis pathology of Cauchy's equation lives only at irrational $p$.
\item Trap recurs when Jensen's inequality is saturated; absent in convex analysis.
\item The same one-line trick proves affineness on convex sets in higher dimensions.
\end{highlights}

\begin{keyword}
Jensen equation \sep Cauchy equation \sep Hamel basis \sep functional equation \sep affine function \sep regularity hypothesis
\MSC[2020] 39B22 \sep 39B05
\end{keyword}

\end{frontmatter}

% =============================================================================
\section{Introduction}
\label{sec:intro}

The equation
\begin{equation}
p\,G(u_1)+(1-p)\,G(u_2) \;=\; G\bigl(p\,u_1+(1-p)\,u_2\bigr)
\qquad (u_1,u_2\in I,\ p\in[0,1]),
\tag{$\star$}\label{eq:star}
\end{equation}
imposed on a function $G\colon I\to\R$ defined on a real interval $I$, has
at least three faces in classical analysis. The faces have been steadily
conflated for over a century, and the conflation has consequences for both
the foundational theory and for the modern applied literature where the
equation surfaces.

\subsection{Three faces of the equation}

The \emph{discrete-coefficient} Jensen equation,
\begin{equation}
G\!\left(\tfrac{u_1+u_2}{2}\right) \;=\; \tfrac{G(u_1)+G(u_2)}{2},
\tag{$J_2$}\label{eq:J2}
\end{equation}
is the $p=\tfrac12$ specialization of~\eqref{eq:star}. Setting
$f(x):=G(x)-G(0)$, the equation~\eqref{eq:J2} is equivalent on a translate
of $I$ to Cauchy's additive equation $f(x+y)=f(x)+f(y)$. Cauchy's equation
inherits a full apparatus of pathological solutions: without measurability,
monotonicity, boundedness on a set of positive Lebesgue measure, or
another regularity hypothesis, the equation admits non-affine solutions
constructed via a \emph{Hamel basis} --- that is, a basis of $\R$ as a
$\Q$-vector space, whose existence requires the axiom of choice. The
classical pathological-solution apparatus is the work of
Hamel~\cite{Hamel1905}; the regularity hypotheses sufficient to recover
affineness were catalogued by Cauchy~\cite{Cauchy1821},
Darboux~\cite{Darboux1875}, Sierpi\'nski~\cite{Sierpinski1920},
Steinhaus~\cite{Steinhaus1920}, and Ostrowski~\cite{Ostrowski1929}.

The \emph{rational-coefficient} Jensen equation,
\begin{equation}
p\,G(u_1)+(1-p)\,G(u_2) \;=\; G\bigl(p\,u_1+(1-p)\,u_2\bigr)
\qquad (u_1,u_2\in I,\ p\in[0,1]\cap\Q),
\tag{$J_{\Q}$}\label{eq:JQ}
\end{equation}
strengthens~\eqref{eq:J2} but does not escape its pathology. Cauchy
additivity entails $\Q$-homogeneity --- $f(qx)=qf(x)$ for every $q\in\Q$
--- by an elementary induction on positive integers, sign reversal for
negatives, and the inverse-of-multiplication argument for $1/n$. Combining
additivity with $\Q$-homogeneity returns~\eqref{eq:JQ} for $f$, hence for
$G$. So \eqref{eq:JQ}, \eqref{eq:J2}, and Cauchy's equation share the same
solution class up to constants, and all three inherit the Hamel-basis
pathology.

The \emph{continuous-coefficient} form~\eqref{eq:star}, by contrast, is
fundamentally different. The strengthening from $\Q$ to $\R$ in the range
of $p$ is \emph{genuine}. As we will see, \eqref{eq:star} does not admit
any Hamel pathology --- the pathology cannot survive at irrational $p$,
and the continuous-coefficient equation tests $G$ at every irrational.
This is the subject of the note.

\subsection{The main result}

\begin{theorem}\label{thm:main}
Let $M>0$ and let $G\colon[0,M]\to\R$ satisfy~\eqref{eq:star} for every
$u_1,u_2\in[0,M]$ and every $p\in[0,1]$. Then $G$ is affine on $[0,M]$ ---
meaning $G(v)=av+b$ for some $a,b\in\R$ --- and explicitly
\[
G(v) \;=\; \frac{G(M)-G(0)}{M}\,v \;+\; G(0).
\]
No regularity hypothesis on $G$ --- continuity, measurability, monotonicity,
boundedness --- is required.
\end{theorem}

The proof is one line: a substitution into~\eqref{eq:star} that pins
$G(v)$ along the chord through the endpoints $(0,G(0))$ and $(M,G(M))$. We
give it in Section~\ref{sec:proof} and use the remainder of the note to
articulate what it means.

\subsection{The historical detour and the structural moral}

Theorem~\ref{thm:main} is folklore (cf.~\cite[\S2.1]{Aczel1966},
\cite[Ch.~13]{Kuczma2009}). It is \emph{not} a new mathematical result.
The contribution of the present note is in three parts: the explicit
dictionary (Section~\ref{sec:dictionary}) of three regularity hypotheses
on $G$ that are required to force affineness for~\eqref{eq:J2} but become
\emph{vestigial} under~\eqref{eq:star}, together with the structural
mechanism (the Hamel pathology lives at irrational $p$, exactly where
\eqref{eq:star}'s continuous coefficient closes the door); the
strict-minimum hypothesis (\S\ref{ssec:weak}) --- the equation
\eqref{eq:star} at a single chord configuration already suffices; and the
recurrence pattern (Section~\ref{sec:recurrence}) of the trap, with three
example application areas (expected-utility theory, Shannon entropy
axiomatic characterization, surrogate calibration on the resolution axis)
where Jensen's inequality is saturated and the one-line substitution of
Section~\ref{sec:proof} is the natural closure.

The 80-year detour through Hamel and the Polish school answered a
perfectly natural question --- \emph{what regularity hypothesis suffices
for \eqref{eq:J2} $\Rightarrow$ affine?} --- but a \emph{different}
question from the one that matters for many modern applications. The
applied derivation does not produce \eqref{eq:J2} or \eqref{eq:JQ}; it
produces \eqref{eq:star} as a saturated identity, and the one-line
substitution of Section~\ref{sec:proof} is the right closure tool. The
dictionary of Section~\ref{sec:dictionary} is the explicit articulation
of which classical regularity hypotheses become vestigial in this
transition.

\paragraph{Terminology.} Throughout the note we refer to the substitution
$u_1 = M$, $u_2 = 0$, $p = v/M$ in~\eqref{eq:star} --- by which $G$ is
pinned along the chord through the endpoints $(0, G(0))$ and $(M, G(M))$
--- as the \emph{chord substitution}. The substitution itself is standard
folklore in the functional-equations literature
(cf.~\cite[\S2.1]{Aczel1966}, \cite[Ch.~13]{Kuczma2009}); the descriptive
label is our convenience, adopted for brevity in repeated reference rather
than as a claim of established terminology.

% =============================================================================
\section{The result and its proof}
\label{sec:proof}

We treat Theorem~\ref{thm:main} in full.

\begin{proof}[Proof of Theorem~\ref{thm:main}]
Fix $v\in[0,M]$. Choose the substitution
\[
u_1 := M, \qquad u_2 := 0, \qquad p := \frac{v}{M}.
\]
The constraints of~\eqref{eq:star} are satisfied: $u_1=M\in[0,M]$,
$u_2=0\in[0,M]$, and $p=v/M\in[0,1]$ since $v\in[0,M]$.

Substituting into~\eqref{eq:star},
\[
\frac{v}{M}\,G(M) + \left(1-\frac{v}{M}\right)G(0)
\;=\; G\!\left(\frac{v}{M}\cdot M + \left(1-\frac{v}{M}\right)\cdot 0\right)
\;=\; G(v).
\]
Rearranging gives the explicit affine formula
$G(v)=G(0)+(G(M)-G(0))\,v/M$.
\end{proof}

The chord substitution pins $G(v)$ along the straight chord from
$(0,G(0))$ to $(M,G(M))$ for every $v\in[0,M]$. Since this determines
$G(v)$ at every $v$, the function $G$ is fully determined by its values at
the two endpoints --- it is exactly the affine function joining them.

\begin{corollary}\label{cor:regularity}
Under the hypotheses of Theorem~\ref{thm:main}, $G$ is in particular
continuous, monotone, locally Lipschitz, absolutely continuous, and
Lebesgue measurable on $[0,M]$.
\end{corollary}

\begin{proof}
Affine functions on bounded intervals have all of these properties: an
affine function $G(v)=av+b$ is differentiable with continuous derivative
$G'(v)=a$, hence continuous; non-decreasing if $a\ge0$ and non-increasing
if $a\le0$; Lipschitz with constant $|a|$; absolutely continuous; and
Borel-measurable, hence Lebesgue-measurable.
\end{proof}

The order of inference matters and is worth stating explicitly. Under
\eqref{eq:star}, the regularity properties of $G$ are the
\emph{conclusion}, not the hypothesis. Each of the regularity hypotheses
one might be tempted to assume --- continuity, monotonicity, boundedness,
measurability --- is automatically satisfied by any solution
of~\eqref{eq:star}. The continuous-coefficient equation supplies its own
regularity from within; no external regularity hypothesis is needed.

This is what makes the continuous-coefficient case structurally different
from the discrete-coefficient case~\eqref{eq:J2}, where some regularity
hypothesis \emph{is} required to rule out the Hamel pathology. We now
examine the difference in detail.

% =============================================================================
\section{The dictionary of vestigial regularity hypotheses}
\label{sec:dictionary}

\subsection{Six hypotheses and their fates}

We place side-by-side the discrete-coefficient equation~\eqref{eq:J2} and
the continuous-coefficient equation~\eqref{eq:star}, and ask which
regularity hypotheses on $G$ are \emph{required} to force affineness for
each. The first column of Table~\ref{tab:dictionary} names the hypothesis
and gives a one-sentence reminder; the second column records the
requirement for~\eqref{eq:J2} with the standard historical attribution;
the third column records what becomes of the requirement
under~\eqref{eq:star}.

\begin{table}[ht]
\renewcommand{\arraystretch}{1.35}
\centering
\begin{tabularx}{\textwidth}{@{} X | X | X @{}}
\toprule
\textbf{Hypothesis on $G$ (definition)} &
\textbf{Required for \eqref{eq:J2}$\Rightarrow$ affine?} &
\textbf{Required for \eqref{eq:star}$\Rightarrow$ affine?}\\
\midrule
\textbf{Continuity on $I$}: $G$ is continuous at every $v\in I$.
& Yes, suffices (Cauchy~\cite{Cauchy1821}).
& \textbf{No} --- continuity is the \emph{conclusion} (Cor.~\ref{cor:regularity}).\\
\midrule
\textbf{Measurability on $I$}: $G$ is Lebesgue- or Borel-measurable.
& Yes, suffices (Sierpi\'nski~\cite{Sierpinski1920}).
& \textbf{No} --- same.\\
\midrule
\textbf{Monotonicity on $I$}: $G$ is non-decreasing or non-increasing.
& Yes, suffices (Darboux~\cite{Darboux1875}).
& \textbf{No} --- same.\\
\midrule
\textbf{Boundedness on a set of positive measure}: there exists
Lebesgue-measurable $E\subseteq I$ with $|E|>0$ on which $G$ is bounded.
& Yes, suffices (Steinhaus~\cite{Steinhaus1920}; cf.~Sierpi\'nski~\cite{Sierpinski1920}).
& \textbf{No} --- same.\\
\midrule
\textbf{Boundedness on $I$}: $G$ is bounded on $I$.
& Yes, suffices (Steinhaus~\cite{Steinhaus1920}; Ostrowski~\cite{Ostrowski1929}).
& \textbf{No} --- same.\\
\midrule
\textbf{None}: no regularity hypothesis at all.
& \textbf{Insufficient} --- Hamel-basis pathology (Hamel~\cite{Hamel1905}).
& \textbf{Sufficient} --- Theorem~\ref{thm:main}.\\
\bottomrule
\end{tabularx}
\caption{Regularity hypotheses on $G$. Column~2 records which hypotheses
are classically required to force affineness for the discrete-coefficient
equation~\eqref{eq:J2}. Column~3 records what they become under the
continuous-coefficient equation~\eqref{eq:star}.}
\label{tab:dictionary}
\end{table}

The bottom row is the point of the table. \textbf{No regularity hypothesis}
is enough for the discrete-coefficient equation~\eqref{eq:J2}: there exists
a non-affine, non-measurable, unbounded, additive function on $\R$ --- the
classical Hamel-basis pathology of Hamel~\cite{Hamel1905} --- that
supplies a non-affine solution of~\eqref{eq:J2}. Conversely, \textbf{no
regularity hypothesis} is needed for the continuous-coefficient
equation~\eqref{eq:star}: Theorem~\ref{thm:main}'s chord substitution
closes the proof at the level of the bare functional equation.

The five non-bottom rows tell a complementary story. For~\eqref{eq:J2},
\emph{some} hypothesis from these five is required; otherwise the Hamel
pathology survives. For~\eqref{eq:star}, \emph{none} of these is required
--- and yet each is automatically satisfied by any solution
(Cor.~\ref{cor:regularity}). The hypothesis is, in the dictionary's
terminology, \emph{vestigial}: classically required for~\eqref{eq:J2},
classically expected for~\eqref{eq:star}, but in fact unnecessary
for~\eqref{eq:star}.

\subsection{The structural mechanism}
\label{ssec:mechanism}

We now examine \emph{why} the Hamel pathology fails to apply
to~\eqref{eq:star}. The mechanism is direct and illuminating.

A Hamel-basis pathological solution $G$ of Cauchy's equation on $\R$ is
\textbf{$\Q$-linear}: by construction (or by induction from additivity),
$G(0)=0$ and $G(qx)=q\,G(x)$ for every $q\in\Q$ and every $x\in\R$. By
design, $G$ is \emph{not} $\R$-linear --- there exists some irrational
$r\in\R$ and some $x\in\R$ such that $G(rx)\ne r\,G(x)$.

The chord identity in~\eqref{eq:star} at the configuration $u_1=M$,
$u_2=0$ reads
\begin{equation}
p\,G(M) + (1-p)\,G(0) \;=\; G(p\,M).
\tag{$\star_0$}\label{eq:star0}
\end{equation}
For a Hamel-pathological $G$ (so $G(0)=0$), \eqref{eq:star0} becomes the
assertion
\[
p\,G(M) \;=\; G(p\,M).
\]
This is exactly the $\R$-linearity assertion for $G$ at the specific pair
$(M,p)$. By $\Q$-linearity of $G$, the assertion holds for every $p\in\Q$.
By the failure of $\R$-linearity, it \emph{fails} at some irrational $p$.
So the Hamel-pathological $G$ satisfies~\eqref{eq:JQ} (every rational
coefficient is OK) but violates~\eqref{eq:star} at every irrational $p$
where the $\R$-linearity defect appears.

The point is structural and worth stating in one sentence.

\begin{quote}
\textbf{The Hamel pathology lives at irrational $p$; the continuous-coefficient equation~\eqref{eq:star} forecloses on it there.}
\end{quote}

The 80-year detour through Hamel and the Polish school --- Cauchy 1821 to
Ostrowski 1929 --- was an investigation of what regularity hypothesis can
substitute for the \emph{unavailable} constraint at irrational $p$, when
only the rational-coefficient version of the equation is in play. With the
continuous-coefficient version, the irrational-coefficient constraint is
\emph{available}, and the entire dictionary collapses.

% =============================================================================
\section{Variants and limits}
\label{sec:variants}

\subsection{The strict-minimum hypothesis}
\label{ssec:weak}

The proof of Theorem~\ref{thm:main} uses~\eqref{eq:star} at exactly one
configuration: $u_1=M$, $u_2=0$, with $p\in[0,1]$ ranging freely. The full
force of~\eqref{eq:star} --- for every pair $(u_1,u_2)$ --- is not
consumed. We record the strict minimum:

\begin{theorem}\label{thm:weak}
Let $M>0$ and let $G\colon[0,M]\to\R$ satisfy
\[
p\,G(M) + (1-p)\,G(0) \;=\; G(p\,M) \qquad \text{for all } p\in[0,1].
\]
Then $G(v)=G(0)+(G(M)-G(0))\,v/M$ on $[0,M]$.
\end{theorem}

The proof is verbatim that of Theorem~\ref{thm:main}: set $p=v/M$.
Theorem~\ref{thm:weak} is \emph{in principle} weaker than
Theorem~\ref{thm:main} --- the hypothesis \eqref{eq:star0} does not a
priori imply the full~\eqref{eq:star}, and as standalone constraints on
$G$ they are inequivalent. Under the conclusion (affineness), both hold;
in practice, an author who derives~\eqref{eq:star} from a richer setup
(e.g., from a Jensen-equality identity in a Bayes-risk computation, see
Section~\ref{sec:recurrence}) typically obtains~\eqref{eq:star} in full
force. The narrower Theorem~\ref{thm:weak} is the right reference for an
author who wants to know how much of the equation the proof actually
consumes.

\subsection{Convex domains in higher dimensions}
\label{ssec:higher}

The chord substitution extends to any convex subset of a real vector space.

\begin{theorem}\label{thm:higher}
Let $V$ be a real vector space, let $C\subseteq V$ be a convex set, and
let $G\colon C\to\R$ satisfy
\[
p\,G(x_1)+(1-p)\,G(x_2) \;=\; G\bigl(p\,x_1+(1-p)\,x_2\bigr)
\qquad (x_1,x_2\in C,\ p\in[0,1]).
\]
Then there is a linear functional $a\colon\mathrm{span}(C-C)\to\R$ and a
constant $b\in\R$ such that $G(x)=a(x-x_0)+b$ on $C$ for any fixed
$x_0\in C$. Equivalently, $G$ is affine on $C$.
\end{theorem}

\begin{proof}[Proof sketch]
Apply Theorem~\ref{thm:main} along every chord in $C$: for each direction
$v$ with $x_0+v\in C$, the map $\lambda\mapsto G(x_0+\lambda v)$ is affine
in the segment parameter $\lambda\in[0,1]$. Hence the increment
$G(x_0+v)-G(x_0)$ depends linearly on $v$ in each direction; an inductive
argument on the dimension of the affine span of the points involved
promotes this directional linearity to a global linear functional $a$ on
$\mathrm{span}(C-C)$. The bookkeeping is standard and is given
in~\cite[Ch.~13]{AczelDhombres1989}.
\end{proof}

The structurally essential step is the one-dimensional chord substitution
of Theorem~\ref{thm:main}. The higher-dimensional lifting is bookkeeping
that transports the one-dimensional affineness into a multidimensional
affineness on the convex set.

\subsection{The rational-coefficient version $(J_{\Q})$ does retain the pathology}
\label{ssec:JQ-pathology}

To sharpen Theorem~\ref{thm:main}'s ``no regularity hypothesis needed'',
we exhibit explicitly that the strengthening from rational to continuous
coefficients in~\eqref{eq:star} is doing real work: the rational-coefficient
version~\eqref{eq:JQ} retains the Hamel pathology.

\begin{proposition}[Folklore]\label{prop:JQ-pathology}
There exists $G\colon[0,1]\to\R$ satisfying~\eqref{eq:JQ}
(hence~\eqref{eq:J2}) on $[0,1]$ that is not affine.
\end{proposition}

\begin{proof}[Construction]
Choose a Hamel basis $H$ of $\R$ over $\Q$ containing $1$. Such a basis
exists by Zorn's lemma applied to the partially ordered set of
$\Q$-linearly independent subsets of $\R$. Define a $\Q$-linear map
$\ell\colon\R\to\R$ by setting $\ell(1)=0$, picking an irrational $h\in
H\cap(0,1)$ (which exists since $H$ has cardinality $\mathfrak c$ and the
rationals plus $\{1\}$ are at most countable, so $H$ must contain
irrationals throughout $\R$ --- in particular in $(0,1)$), setting
$\ell(h)=1$, and extending arbitrarily on the remaining basis elements.

By construction, $\ell$ is $\Q$-linear: $\ell(x+y)=\ell(x)+\ell(y)$ for
all $x,y\in\R$ (so $\ell$ satisfies Cauchy's equation) and
$\ell(qx)=q\ell(x)$ for all $q\in\Q$ and $x\in\R$ ($\Q$-homogeneity). But
$\ell$ is \emph{not} $\R$-linear: $\ell(1)=0$ while $\ell(h)=1$, so the
values of $\ell$ on the $\Q$-span of $\{1\}$ (namely $\Q$, on which $\ell$
vanishes identically) and on the $\Q$-span of $\{h\}$ (where $\ell$ is
nonzero) are incompatible with any $\R$-linear functional on $\R$.

Set $G:=\ell|_{[0,1]}$. Then $G$ satisfies~\eqref{eq:JQ} on $[0,1]$
because $\Q$-rational convex combinations of points of $[0,1]$ remain in
$[0,1]$ and $\ell$ respects $\Q$-linear combinations on all of $\R$.

Explicit witness that $G$ is not affine: $G(h)=\ell(h)=1\ne 0$, while
$G(q)=\ell(q)=q\,\ell(1)=0$ for every rational $q\in[0,1]$. Any affine map
$A\colon[0,1]\to\R$ agreeing with $G$ on $\Q\cap[0,1]$ would satisfy
$A(0)=0$ and $A(q)=0$ for all rational $q$ --- forcing $A\equiv 0$,
contradicting $G(h)=1$. So $G$ is not affine.
\end{proof}

Proposition~\ref{prop:JQ-pathology} affirms the dictionary: \eqref{eq:JQ}
on $[0,1]$ admits non-affine solutions; only the strengthening to the
\emph{continuous}-coefficient~\eqref{eq:star} retires the
regularity-hypothesis requirement.

% =============================================================================
\section{Where the trap recurs: Jensen saturation as a structural pattern}
\label{sec:recurrence}

The chord substitution of Section~\ref{sec:proof} is, as noted, folkloric
--- \cite[\S2.1]{Aczel1966} and \cite[Ch.~13]{Kuczma2009} articulate the
underlying observation, and the result is standard in the
functional-equations community. Why, then, does this note exist?

Because the \emph{vestigial regularity hypotheses} --- continuity,
measurability, boundedness --- keep reappearing in applied derivations
of~\eqref{eq:star} outside the functional-equations community. The trap
is not a one-off oversight; it has a structural source. This section
identifies the source and surveys three example application areas where
it surfaces.

\subsection{The structural source: saturated Jensen}
\label{ssec:saturated-jensen}

The recurrence is predictable. Whenever a derivation in any field arrives
at an identity of the form
\begin{equation}
\E[g(\xi)] \;=\; g\bigl(\E[\xi]\bigr)
\qquad (\xi \text{ a random variable on } I,\ g\colon I\to\R),
\tag{$\dagger$}\label{eq:dagger}
\end{equation}
\emph{for a class of random variables $\xi$ wide enough that the marginal
$\E[\xi]$ can be any point of $I$ and the support of $\xi$ can be any
two-point subset of $I$ with any pair of masses $(p,1-p)$ with $p\in[0,1]$},
the identity is exactly~\eqref{eq:star} with $g=G$ and $\xi$ supported on
$\{u_1,u_2\}$ with respective masses $\{p,1-p\}$. Such an identity arises
whenever \textbf{Jensen's inequality is pushed to saturation} --- that is,
whenever equality is required in $\E[g(\xi)]\le g(\E[\xi])$ for concave
$g$ across the full two-point family.

Jensen's inequality is the central tool in dozens of fields: convex
analysis, information theory, statistical decision theory, mathematical
economics, mathematical physics. In \emph{most} of these fields, Jensen is
used as an inequality with slack --- the slack is explicitly tracked
(in convex analysis via biconjugates and Fenchel duality; in calibration
theory via the $\psi$-transform construction of Bartlett, Jordan, and
McAuliffe~\cite{BJM2006}; in entropy via convexity arguments). Such
derivations never produce~\eqref{eq:star} as a saturated identity.

In \emph{some} fields, however, a derivation produces~\eqref{eq:dagger} as
a saturated identity --- the slack is forced to zero by a structural
property of the setup. At that moment~\eqref{eq:star} appears, and the
chord substitution is the natural closure. An author unfamiliar with the
chord substitution will reach for the classical Cauchy-Hamel toolkit and
find that the regularity hypothesis is unnecessary only after the fact.

\subsection{Application: expected-utility representation theorems}
\label{ssec:utility}

The first application area is \textbf{expected-utility theory} in the
von Neumann--Morgenstern tradition~\cite{vNM1944}. The 1944 axiomatization
of preferences over lotteries, sharpened by Herstein and Milnor in
1953~\cite{HersteinMilnor1953}, derives a utility functional $U$ on a
space of lotteries $\mathcal L$ that is \emph{linear in probability}:
\[
U\bigl(p\,L_1+(1-p)\,L_2\bigr) \;=\; p\,U(L_1)+(1-p)\,U(L_2)
\qquad (L_1,L_2\in\mathcal L,\ p\in[0,1]).
\]
This is exactly~\eqref{eq:star} with the abstract lotteries $L_1,L_2$ as
the points $u_1,u_2$ and $U$ as the function $G$. The classical closure of
the linearity step in the Herstein--Milnor proof uses the
\emph{Archimedean axiom}: for any three lotteries $L_1\succ L_2\succ L_3$,
there exist $p,q\in(0,1)$ such that
$pL_1+(1-p)L_3\succ L_2\succ qL_1+(1-q)L_3$. The Archimedean axiom forces
the continuity-in-probability of preferences, which in turn forces
continuity-in-probability of $U$.

But the chord substitution of Theorem~\ref{thm:main} \emph{also} closes
the linearity step, without requiring any continuity hypothesis. Both
routes work; the chord substitution is the algebraic alternative to the
Archimedean axiom. Even where the Archimedean axiom is invoked (and
rightly so, as it carries other preference-theoretic content), the
closure of the linearity equation itself does not require it.

\subsection{Application: Shannon entropy's axiomatic characterization}
\label{ssec:entropy}

The second application area is the \textbf{axiomatic characterization of
Shannon entropy} in the tradition of Khinchin~\cite{Khinchin1957} and
Faddeev~\cite{Faddeev1957}. Shannon's 1948 paper introduced the entropy
function $H(p_1,\ldots,p_n)=-\sum_i p_i\log p_i$; the question of which
conditions on a continuous function on $\bigcup_n\Delta^{n-1}$ uniquely
determine $H$ up to a multiplicative constant was answered by Khinchin and
Faddeev. Their axioms are: continuity in the arguments; maximum at the
uniform distribution; and a \emph{recursivity axiom} (Khinchin--Faddeev)
that relates the entropy of a joint distribution to a marginal entropy
plus a conditional.

In the modern formalization (\cite[Ch.~5]{Aczel1966},
\cite[\S22]{AczelDhombres1989}), the recursivity axiom reduces, in
intermediate steps, to a system of functional equations on auxiliary
functions $g,h$ on $[0,1]$. Several of these intermediate equations are
of \eqref{eq:star}-form for the auxiliary functions, and the closure step
in the modern axiomatic treatment uses continuity (the first Khinchin
axiom). The chord substitution of Theorem~\ref{thm:main} is the
alternative closure that does not require continuity-as-hypothesis.

\subsection{Application: surrogate calibration on the resolution axis}
\label{ssec:calibration}

The third application area is \textbf{surrogate calibration} in modern
statistical decision theory. The foundational tradition is Bartlett,
Jordan, and McAuliffe~\cite{BJM2006} (binary classification) and its
multiclass extensions: Tewari and Bartlett~\cite{TewariBartlett2007},
Steinwart~\cite{Steinwart2007}, Reid and
Williamson~\cite{ReidWilliamson2010,ReidWilliamson2011}. These works study
the question of how minimizing a convex surrogate loss approximately
minimizes the 0--1 loss in classification.

The standard derivation route in the surrogate-calibration literature uses
\emph{convex analysis} --- biconjugates, supporting hyperplanes, Fenchel
duality, the $\psi$-transform of Bartlett--Jordan--McAuliffe --- and
consequently tracks Jensen's inequality as an \emph{inequality} with
explicit slack. The slack is the calibration function $\psi$, and the
construction never produces~\eqref{eq:star} as a saturated identity.

In a recent manuscript of the author~\cite{El2} on the achievable error
floor of partition-based classifiers, the resolution-axis transposition
framing of surrogate calibration produces~\eqref{eq:star} directly: a
two-cell partition computation forces~\eqref{eq:star} on the function $G$
that expresses the partition Bayes risk in terms of an aggregated concave
score, with the cell mass $p$ ranging freely over $[0,1]$ on an atomless
underlying probability space. In the Lean~4 formalization of~\cite{El2},
the corresponding lemma was initially declared with a boundedness
hypothesis on $G$ in deference to the Cauchy-Hamel literature; the proof
body then exhibited that the hypothesis was unused, by the chord
substitution of Theorem~\ref{thm:main}.

The lesson generalizes beyond~\cite{El2}. Any derivation in any field that
saturates Jensen across a wide class of two-point distributions should be
expected to produce~\eqref{eq:star} similarly. At that moment,
Theorem~\ref{thm:main} closes the regularity-hypothesis question on the
spot.

\subsection{An invitation}
\label{ssec:invitation}

The three application areas of \S\S\ref{ssec:utility}--\ref{ssec:calibration}
are not exhaustive. We invite readers who have encountered~\eqref{eq:star}
in their own work (in mathematical economics, in information theory, in
statistical decision theory, in physics, or elsewhere) and who have
invoked an unnecessary Cauchy-Hamel regularity hypothesis to extend the
catalog. The structural pattern of saturated Jensen --- and the chord
substitution as its natural closure --- is, we believe, recurrent enough
to merit explicit articulation.

% =============================================================================
\section{Concluding remarks}
\label{sec:conclusion}

Theorem~\ref{thm:main} is, mathematically, one line. The dictionary of
Section~\ref{sec:dictionary} is six rows of a table. The structural
argument of \S\ref{ssec:mechanism} is one sentence: \emph{the Hamel
pathology lives at irrational $p$; the continuous-coefficient
equation~\eqref{eq:star} forecloses on it there.} Why does this material
deserve a Short Research Note?

Two reasons. First, the 80-year detour through Hamel and the Polish school
is a \emph{foundational} result of classical analysis --- and the
dictionary of Section~\ref{sec:dictionary} articulates a precise sense in
which the detour answered an adjacent question: not ``what is the minimum
hypothesis for~\eqref{eq:star} $\Rightarrow$ affine?''
(Theorem~\ref{thm:main}'s answer: nothing), but ``what is the minimum
hypothesis for~\eqref{eq:J2} $\Rightarrow$ affine?'' (Polish school's
answer: any of a dictionary of tameness conditions). Both questions are
interesting; their answers are different.

Second, the recurrence pattern of Section~\ref{sec:recurrence} is, to our
knowledge, not articulated explicitly elsewhere in the functional-equations
literature. The pattern --- saturated Jensen produces~\eqref{eq:star},
and the chord substitution closes it --- is recognizable in a number of
modern application contexts (utility theory, entropy, calibration), and
the explicit articulation may save other authors from carrying vestigial
regularity hypotheses in their own work.

The chord substitution itself is the kind of one-line proof that, once
seen, is hard to unsee --- but that nonetheless takes some scaffolding to
make visible. The present note is the scaffolding.

% =============================================================================
\section*{Funding}

This research did not receive any specific grant from funding agencies in
the public, commercial, or not-for-profit sectors.

\section*{Acknowledgements}

The chord-substitution observation surfaced in the Lean~4 formalization
track of the author's main paper~\cite{El2}, during a development phase
whose discipline is captured in three open project skills documented in
the source code repository associated with~\cite{El2}. The author thanks
the internal adversarial-review process described in the same source for
surfacing the over-engineered boundedness hypothesis whose retirement
motivated the present note.

\section*{Declaration of competing interests}

The author declares that they have no known competing financial interests
or personal relationships that could have appeared to influence the work
reported in this paper. The corresponding competing-interests declaration
is provided in a separate file accompanying this submission.

\section*{Declaration of generative AI and AI-assisted technologies in the manuscript preparation process}

During the preparation of this work the author used Anthropic's Claude
Opus (accessed via GitHub Copilot) in two specific capacities:
(i)~\emph{adversarial review} --- the multi-round audit cycle documented
in the source repository accompanying this work, in which the model was
prompted in hostile-referee mode to surface defects in the prose, the
argument structure, and the prior-art positioning of each draft;
(ii)~\emph{paper-formatting tasks} --- LaTeX mechanics (Elsevier
\texttt{elsarticle} class configuration, theorem-environment setup, table
layout), BibTeX entry construction, cross-reference checking, and
Markdown-to-LaTeX conversion. The mathematical content --- including all
theorems, proofs, examples, the structural argument of
Section~\ref{sec:recurrence}, and the dictionary of regularity hypotheses
in Section~\ref{sec:dictionary} --- was conceived, verified, and refined
by the author. The author has reviewed and edited all content as needed
and takes full responsibility for the published article.

% =============================================================================
% Bibliography (Elsevier numerical style, references in citation order)
\bibliographystyle{elsarticle-num}
\bibliography{refs}

\end{document}
